\documentclass[10pt,a4paper]{amsart}
\usepackage[margin=19mm]{geometry}
\usepackage{amsmath,amssymb,mathtools}
\usepackage[scr=rsfs,cal=euler]{mathalfa}
\usepackage{xfrac}
\usepackage[unicode,hidelinks,bookmarks=false]{hyperref}
\newcommand{\ie}{i.e.\ }
\newcommand{\cf}{cf.\ }
\newcommand{\resp}{respectively\ }
\newcommand{\Subsection}{\subsection}
\providecommand{\nobreakdash}{\nobreak\hbox{-}}
\setlength{\emergencystretch}{2em}
\multlinegap0pt
\usepackage{booktabs}
\usepackage{nicefrac}
\usepackage{xcolor}
\usepackage{float}
\usepackage{bm}
\usepackage{amssymb}
\usepackage{tikz}
\newtheorem{lemma}{Lemma}
\title{On well-posed energy/entropy stable boundary conditions for the rotating shallow water equations}
\author{Kenneth Duru, Chuqiao Xu}
\date{Source: arXiv:2601.02513v1 (CC BY 4.0)}
\begin{document}
\maketitle

\noindent\emph{Source and licence.} On well-posed energy/entropy stable boundary conditions for the rotating shallow water equations. Version arXiv:2601.02513v1 (CC BY 4.0), distributed by arXiv under the Creative Commons Attribution 4.0 International licence, http://creativecommons.org/licenses/by/4.0/, as recorded in the arXiv licence metadata for this exact version and shown on the abstract page https://arxiv.org/abs/2601.02513v1. Full licence notice: \texttt{licenses/arxiv-2601.02513-CC-BY-4.0.txt}.\par\smallskip

\noindent\textit{Editorial scope.} Counted unit: the article's lemma lem:boundary-term-pointwise-nonlinear (source0.tex lines 1528--1537). The article's complete original proof of it is delivered with it (source0.tex lines 1538--1605, one \texttt{proof} environment of 68 lines). No mathematics is written, repaired or re-derived: the statement and the proof are the article's own lines, in the article's own notation, and no retained line is substituted or re-flowed.  Retained uncounted as the source-local context the proof needs: lines 1510--1519.  Nothing was omitted from any retained span, so the package carries no omission record.  No retained line contains a citation, so the wrapper carries no bibliography and no bibliography entry.  No retained line contains a figure environment, an image-inclusion command or drawing code, so no image is delivered and the package declares no asset dependency.  Editorial formatting only, in addition: an AMS \texttt{amsart} preamble that restates the article's own declarations and macros used by the retained lines, namely the environments \texttt{lemma} on separate counters, together with the standard packages those lines need.  The retained environments print in this document as Lemma 1 in order of appearance, whereas the article prints the same items with the numbering of its own full document.  The complete native source of the article as distributed in the arXiv e-print for this exact version is bundled unchanged as \texttt{sources/192154/source0.tex}.
\begin{equation}\label{eq:boundary-term-pointwise-nonlinear}
    \operatorname{\mathbb{BT}}_j =
\begin{cases}
-\left(G_{j}F_{nj}\right)+ \tau_h G_{j} \left(\alpha G_{nj} - \beta F_{nj}\right) +
     \tau_n F_{nj}\left(\alpha G_{nj} - \beta F_{nj}\right) + \tau_s h_{j} u_{nj} u_{sj}^2, \quad \ \text{if} \ u_{nj} < 0\\
%     
-\left(G_{j}F_{nj}\right)+ \tau_h G_{j} \left(\alpha G_{nj} - \beta F_{nj}\right) +
     \tau_n F_{nj}\left(\alpha G_{nj} - \beta F_{nj}\right) , \quad \ \text{if} \ u_{nj} \ge  0.
     \end{cases}.
\end{equation}

\begin{lemma}\label{lem:boundary-term-pointwise-nonlinear}
    Consider the boundary term  $\operatorname{\mathbb{BT}}_j$ defined by \eqref{eq:boundary-term-pointwise-nonlinear}, where $u_n + c >0$, $u_n-c < 0$. and assume $\alpha\ge 0$,  $\beta \ge 0$ with $|\alpha| + |\beta| >0$.
    If the boundary parameters $\alpha$, $\beta$ and the penalty weights $\tau_h$, $\tau_n$ and $\tau_s$ satisfy
    \begin{itemize}
      \item[ 1:]  When $\alpha = 0$, $\beta > 0$:  $\tau_h = -1/\beta$, $\tau_n \ge 0$ and $\tau_s \ge 1/2$.
       \item[2:]  When  $\alpha > 0$, $\beta = 0$:  $\tau_h \le 0$, $\tau_n = 1/\alpha$ and $\tau_s \ge 1/2$,
       \item[3:]  When $\alpha > 0$, $\beta > 0$:  $\tau_h =-1/(2\beta)$, $\tau_n = 1/(2\alpha)$ and $\tau_s \ge 1/2$,
    \end{itemize}
  then the boundary is never positive, that is $\operatorname{\mathbb{BT}}_j \le 0$.
\end{lemma}

\begin{proof}
Case 1:
    Consider the boundary term  $\operatorname{\mathbb{BT}}_j$ defined by \eqref{eq:boundary-term-pointwise-nonlinear} with $\alpha= 0$ and $\beta >0 $, and set $\tau_h = -1/\beta$
     we have
     $$
      \operatorname{\mathbb{BT}}_j =
\begin{cases}
     -\tau_n \beta F_{nj}^2 + \tau_s h_j u_{nj} u_{sj}^2, \quad \ \text{if} \ u_{nj} < 0\\
%     
     -\tau_n \beta F_{nj}^2  , \quad \ \text{if} \ u_{nj} \ge  0.
     \end{cases}.
     $$
     Thus if $\tau_n \ge 0$ and $\tau_s \ge 1/2$ then $\operatorname{\mathbb{BT}}_j\le 0$.
%   \begin{equation*}
%     \operatorname{\mathbb{BT}}_j =
% \begin{cases}
% -\left(\lambda_1 + \lambda_2 \gamma^2\right)w_{1j}^2 +  U_n (\tau_s-1) w_{3j}^2, \quad \ \text{if} \ U_n < 0\\
     
% -\left(\lambda_1 + \lambda_2 \gamma^2\right)w_{1j}^2  - U_n w_{3j}^2
%       , \quad \ \text{if} \ U_n \ge  0.
%      \end{cases}
% \end{equation*} 

\noindent
Case 2:
    Consider the boundary term  $\operatorname{\mathbb{BT}}_j$ defined by \eqref{eq:boundary-term-pointwise-nonlinear} with $\alpha> 0$ and $\beta =0 $, and set $\tau_n = 1/\alpha$
     we have
     $$
      \operatorname{\mathbb{BT}}_j =
\begin{cases}
     \tau_h \alpha G_{j}  G_{nj}  + 
    \left(\tau_s -\frac{1}{2}\right)h_ju_{nj}u_{sj}^2 , \quad \ \text{if} \ u_{nj} < 0\\
%     
 \tau_h \alpha G_{j}  G_{nj}  
     -\frac{1}{2}h_ju_{nj}u_{sj}^2, \quad \ \text{if} \ u_{nj} \ge  0.
     \end{cases}.
     $$
where we have used
$
G_n = gh + \frac{1}{2}hu_n^2 
$
and 
$G = gh + \frac{1}{2}u_n^2 + \frac{1}{2}u_s^2  = G_n + \frac{1}{2}u_s^2 $.
Note that $G_{j} >0$,  $G_{nj} >0$, and $G_{j}  G_{nj} >0 $, so if $\alpha >0$, $\tau_h \le 0$ and $\tau_s \ge 1/2$ then $\operatorname{\mathbb{BT}}_j \le 0$.
\newline
\newline
\noindent
Case 3:
    Consider the boundary term  $\operatorname{\mathbb{BT}}_j$ defined by \eqref{eq:boundary-term-pointwise-nonlinear} with $\alpha> 0$ and $\beta >0 $, and set $\tau_h = -1/(2\beta)$ $\tau_n = 1/(2\alpha)$
     we have
     $$
      \operatorname{\mathbb{BT}}_j =
\begin{cases}
      -\frac{\alpha}{2\beta} G_{j}  G_{nj}  - \frac{\beta}{2\alpha} F_{nj}^2 + 
    \left(\tau_s -\frac{1}{4}\right)h_ju_{nj}u_{sj}^2 , \quad \ \text{if} \ u_{nj} < 0\\
%     
-\frac{\alpha}{2\beta} G_{j}  G_{nj}  - \frac{\beta}{2\alpha} F_{nj}^2  
     -\frac{1}{4}h_ju_{nj}u_{sj}^2, \quad \ \text{if} \ u_{nj} \ge  0.
     \end{cases}.
     $$
As above, we have used
$
G_n = gh + \frac{1}{2}u_n^2 
$
and 
$G = gh + \frac{1}{2}u_n^2 + \frac{1}{2}u_s^2  = G_n + \frac{1}{2}u_s^2 $.
Note that $G_{j} >0$,  $G_{nj} >0$, and $G_{j}  G_{nj} >0 $, thus if  $\tau_s \ge 1/2 > 1/4$ then $\operatorname{\mathbb{BT}}_j \le 0$.
\end{proof}

\end{document}
